\section{Nested support bounds from exterior minors}

\subsection*{Finding F-016: tail-sum support and full determinant multiplicity}
Let \(\lambda\vdash nd\), with
\(r=\ell(\lambda)\leq d\), and let \(U\) have basis
\(e_1,\ldots,e_r\). Every highest-weight vector of weight
\(\lambda\) in \(\operatorname{Sym}^d\operatorname{Sym}^n U\)
belongs to \(\operatorname{Sym}^d W_{n,\lambda}\), where
\begin{equation}\label{eq:F016support}
 W_{n,\lambda}=
 \operatorname{span}\left\{e^\eta:|\eta|=n,\quad
             \sum_{i=j}^r\eta_i\leq\lambda_j
                             \text{ for every }2\leq j\leq r\right\}.
\end{equation}
This assertion does not require \(n\geq\lambda_2\).
It constrains the sum of the exponents in each terminal set
of rows, and is stronger than the individual coordinate bounds.

\begin{proof}
Use the matrix polynomial model
\(\operatorname{Sym}(U\otimes\mathbb C^d)\), with row
variables indexed by \(i\) and column variables by \(h\).
As proved in F-014, the entire row highest-weight space of
weight \(\lambda\) is spanned by products of flag minors
whose heights are the column lengths of \(\lambda\).
For completeness, their span is a nonzero column-group-stable
subspace of the irreducible Cauchy module
\(S_\lambda\mathbb C^d\), and hence is the whole space.

Exactly \(\lambda_j\) minor factors in each such product
have height at least \(j\). A minor of smaller height uses
no row \(j,\ldots,r\). In a monomial of a minor of height
at least \(j\), any fixed column contributes at most one
variable. Consequently every monomial in every spanning product,
and therefore every polynomial in the entire highest-weight
space, has degree at most \(\lambda_j\) in rows
\(j,\ldots,r\) of any one fixed column.

Now take the column multidegree \((n,\ldots,n)\) subspace.
Its monomial basis is naturally the tensor-product monomial
basis of \((\operatorname{Sym}^n U)^{\otimes d}\).
The preceding bound holds separately in every column, so
the relevant highest-weight subspace lies in
\(W_{n,\lambda}^{\otimes d}\). Taking the column-permutation
invariants gives
\[
       (W_{n,\lambda}^{\otimes d})^{S_d}
                                  =\operatorname{Sym}^d W_{n,\lambda}.
\]
Under this identification, an unordered product of monomial
generators is the sum of its distinct column assignments,
up to a nonzero scalar. Different unordered products have
disjoint sets of assignments. Thus no cancellation in this
passage could conceal a forbidden coefficient generator.
This proves \eqref{eq:F016support} for every copy.
\end{proof}

Suppose now \(r\geq2\), put \(b=\lambda_2\), and define
the finite, coordinatewise downward-closed set
\[
 \Gamma_\lambda=
 \left\{\alpha\in\mathbb Z_{\geq0}^{r-1}:
          \sum_{i=j}^r\alpha_i\leq\lambda_j
                              \text{ for all }2\leq j\leq r\right\}.
\]
In particular \(|\alpha|\leq b\). For \(s\geq0\), write
\[
 C_s^{\mathrm{flag}}(\lambda)=
       \#\{\alpha\in\Gamma_\lambda:|\alpha|\leq s\}.
\]
For a subset \(A\subseteq\{2,\ldots,r\}\), write
\(\alpha|_{A^c}=0\) to mean that its coordinates outside
\(A\) vanish, and put
\[
 \begin{split}
 N_A^{\mathrm{flag}}(\lambda)&=
       \#\{\alpha\in\Gamma_\lambda:\alpha|_{A^c}=0\},\\
 D_{\mathrm{flag}}(\lambda)&=
 C_{\lfloor(b-1)/2\rfloor}^{\mathrm{flag}}(\lambda)
 +C_{\lceil(b-1)/2\rceil}^{\mathrm{flag}}(\lambda)-1,\\
 T_{\mathrm{flag}}(\lambda)&=
 \min_{A\sqcup B=\{2,\ldots,r\}}
       \bigl(N_A^{\mathrm{flag}}(\lambda)
                         +N_B^{\mathrm{flag}}(\lambda)-1\bigr).
 \end{split}
\]
If also \(r\leq n^2\), then
\begin{equation}\label{eq:F016K}
 n\geq\min\{D_{\mathrm{flag}}(\lambda),T_{\mathrm{flag}}(\lambda)\}
 \quad\Longrightarrow\quad
                     K_n(\lambda)=a_\lambda(d[n]).
\end{equation}

\begin{proof}
All powers of the second variable up to degree \(b\) lie
in \(\Gamma_\lambda\). Exactly as in F-015,
\(D_{\mathrm{flag}}\geq b\) and
\(T_{\mathrm{flag}}\geq b+1\). Thus \(n\geq b\).
Apply F-014 to obtain
\[
 \mu=(db-|\bar\lambda|,\lambda_2,\ldots,\lambda_r),
                 \qquad a_\mu(d[b])=a_\lambda(d[n]).
\]
The first part of this finding, now applied at degree \(b\)
and weight \(\mu\), puts its entire highest-weight space
in \(\operatorname{Sym}^d W\), where
\[
 W=\operatorname{span}
      \{e_1^{b-|\alpha|}e_2^{\alpha_2}\cdots e_r^{\alpha_r}:
                                        \alpha\in\Gamma_\lambda\}.
\]
The tail bounds for \(\mu\) and \(\lambda\) are identical.

Let \(\mathcal C\) be the corresponding dual monomial
coefficient space of degree-\(b\) forms. The graph construction
O-015, or the constructions O-014.a and M-009 restricted to
these downward-closed vertex sets, puts
\(x_1^{n-b}\mathcal C\) in \(\Omega_n\) under the respective
degree-split or variable-split inequality. Indeed, every
coordinatewise smaller multi-index satisfies all terminal-sum
bounds. Splitting a monomial between the trees, or first
removing the variable carried by its coefficient edge, keeps
both endpoints and every tree path in the permitted sets.
The cycle-cover signs and exponents are unchanged.

The restriction-and-padding proof of F-015 applies to this
monomial space \(W\) without any assumption that it is
Borel-stable. Restriction to \(\mathcal C\) is injective
on \(\operatorname{Sym}^d W\), and the composition of the
highest-weight lift with actual padding pullback acts on each
generator by the nonzero scalar
\[
                       \frac{(b-|\alpha|+n-b)!}
                            {(b-|\alpha|)!}.
\]
Thus all \(a_\mu(d[b])=a_\lambda(d[n])\) copies remain
independent in the determinant coordinate ring. The ambient
upper bound gives \eqref{eq:F016K}.
\end{proof}

\paragraph{A simple bound from a cut in the rows.}
For \(2\leq k\leq r\), split the variables into rows
\(2,\ldots,k-1\) and \(k,\ldots,r\). Ignoring other
terminal-sum bounds can only increase the counts. The first
part has \(k-2\) variables of total degree at most \(b\),
and the second has \(r-k+1\) variables of total degree at
most \(\lambda_k\). Hence
\begin{equation}\label{eq:F016cut}
 T_{\mathrm{flag}}(\lambda)
 \leq \binom{b+k-2}{k-2}
           +\binom{r-k+1+\lambda_k}{\lambda_k}-1.
\end{equation}
If the right-hand side is at most \(n\), every ambient
copy of \(\lambda\) occurs on the determinant closure.
At \(k=3\), this reads
\begin{equation}\label{eq:F016thirdrow}
 n\geq\lambda_2+\binom{r-2+\lambda_3}{\lambda_3}
             \quad\Longrightarrow\quad K_n(\lambda)=a_\lambda(d[n]).
\end{equation}
In particular, for \(r\geq3\) and \(\lambda_3\leq2\),
the readily checked sufficient condition is
\[
                       n\geq\lambda_2+\binom r2.
\]
This allows second rows almost as large as \(n\), even at
arbitrarily large outer degree. It is much stronger on these
tails than the uniform constant second-row exclusion alone.

For the exact tail \((b,2^{r-2})\), \(b\geq2\), the nested
condition is simply that the variables after the second have
total degree at most two. Its degree-split counts, for \(s\geq1\),
are
\[
 C_s^{\mathrm{flag}}
  =(s+1)+s(r-2)+(s-1)\binom{r-1}{2}
  =r+(s-1)\binom r2,
\]
where the displayed formula is used only for \(s\leq b\).
Thus \(D_{\mathrm{flag}}=2r+(b-3)\binom r2-1\) for
\(b\geq3\), and \(D_{\mathrm{flag}}=r\) for \(b=2\).
The variable-split bound \(b+\binom r2\) can be smaller still.

\paragraph{Relation to the actual GCT comparison.}
The support statement \eqref{eq:F016support} applies directly
to determinant-ideal highest-weight functions before any
stabilization. After aligning the padding form, it justifies
deleting degree-\(m\) coefficients outside the same terminal-sum
constraints while preserving evaluation. The root finding
M-011 develops that variant without requiring \(n\geq\lambda_2\).
By contrast, \eqref{eq:F016K} proves the stronger equality
with the full ambient multiplicity, using the stabilized base
degree \(b=\lambda_2\). These two consequences have different
hypotheses and neither is inferred from occurrence alone.

\paragraph{Grothendieck dictionary.}
The additional information comes from exterior-power matrix
coefficients: a flag minor cannot use the same column twice.
The Cauchy decomposition identifies their column-group span
with the complete Schur multiplicity module, as in Macdonald,
I(4.3), p.~63, and Appendix~A(5.3)--(5.6), pp.~155--156.
Passing to a terminal row set is a support calculation in the
torus grading of the matrix scheme, under SGA~3, I, 4.7.3.
The permitted coefficient support defines an affine linear
subscheme; it need only be stable under the diagonal torus.
The associated determinant morphisms are integral identities.
The all-copy conclusion uses characteristic-zero highest-weight
modules and the factorial pairing, with exactly the same
base-change qualifications as F-015. It supplies an explicit
condition under which the actual coherent quotient loses no
copies of the indicated Schur functor.
